\documentclass[11pt]{article}
\textwidth 15.9cm
\textheight 23. cm
\addtolength{\oddsidemargin}{-11.5mm}
\addtolength{\evensidemargin}{-11.5mm}
\addtolength{\topmargin}{-19.0mm}

\usepackage[ansinew]{inputenc}
\usepackage{amsmath,amssymb}
\usepackage{bbm}
\usepackage{graphicx}
\usepackage{hyperref}

\include{GempicLatexMacros}

\renewcommand{\part}{\texttt{part}}
\newcommand{\cell}{\texttt{cell}}
\newcommand{\node}{\texttt{node}}
\newcommand{\prim}{\texttt{prim}}

\begin{document}

We currently have two types of solvers: Poisson solvers and a quasineutral solver.
\section{Poisson Solvers}
The Poisson solvers are used for solving (mild variations of) the Poisson equation,
\begin{equation}
    -\Delta\phi=\rho.
\end{equation}

The individual solvers are not supposed to be used directly, but rather selected from the input file as detailed in \ref{sec:PoissonSolverInput} and called in the code using the pointer generated by calling \lstinline{Gempic::FieldSolvers::make_poisson_solver}.

\subsection{Amrex}
The \lstinline{AMReX} solver uses the \lstinline{amrex::MLMG} class with \lstinline{MLNodeLaplacian} as the linear operator (cf. the \lstinline{AMReX}~\href{https://amrex-codes.github.io/amrex/docs_html/LinearSolvers.html}{linear solver documentation}) to solve the nodal equation
\begin{equation}
    \tildearr{\rho}=-\Delta\tildearr{\phi}
\end{equation}
for $\tildearr{\phi}$, which we subsequently convert to $\arr{\phi}$. Requires periodic boundary conditions and enforces that the average value of $\tildearr{\rho}$ is $0$.

\subsection{ConjugateGradient}
The conjugate gradient solver solves the equation 
\begin{equation}
    \tildearr{\rho}=-\tildemat{D}\mat{H}_2\mat{G}\arr{\phi}
\end{equation}
for $\arr{\phi}$ using the conjugate gradient method. With periodic boundary conditions, this solver enforces that the average value of $\tildearr{\rho}$ is $0$. The order of the solver is equal to the Hodge degree.

\subsection{ConjugateGradientInverseHodge}
The conjugate gradient inverse Hodge solver solves the equation 
\begin{equation}
    \tildearr{\rho}=-\tildemat{D}\mat{H}_{2,CG}\mat{G}\arr{\phi}
\end{equation}
for $\arr{\phi}$ using the conjugate gradient method.
Here, $\mat{H}_{2,CG}$ is itself a conjugate solver step solving $\tildemat{H}_2(\tilde{\mat{G}\arr{\phi}})=(\mat{G}\arr{\phi})$ for $(\tilde{\mat{G}\arr{\phi}})$.
This is an option since the Hodge is generally not its own inverse, \emph{i.e.}, $\tildemat{H}\neq\mat{H}^{-1}$, but $\mat{H}_{2,CG}=\mat{H}^{-1}$.
With periodic boundary conditions, this solver enforces that the average value of $\tildearr{\rho}$ is $0$. The order of the solver is equal to the Hodge degree - 2 (and 2 for Hodge degree 2, for which the Hodge is diagonal).

\subsection{FFT}
The FFT solver solves the equation 
\begin{equation}
    \tildearr{\rho}=-\tildemat{D}\mat{H}_2\mat{G}\arr{\phi}
\end{equation}
for $\phi$ using the FFT forward ($\mathcal{F}$) and -backward operators ($\mathcal{F}^{-1}$).
The solution is found as
\begin{equation}
    \arr{\phi} = \mathcal{F}^{-1}(\mathcal{F}(\tildearr{\rho})/\lambda_{total})
\end{equation}
where $\lambda_{total} = \lambda_x + \lambda_y + \lambda_z$ are the eigenvalues of the diagonal form of $\tildemat{D}\mat{H}_2\mat{G}$ and
\begin{equation}
    \lambda_{i_k} = (2-2\cos(2\pi k/N_x))\lambda^k_{\mat{H}_2}
\end{equation}
with $N_x$ the number of grid points and $\lambda^k_{\mat{H}_2}$ the $k^{th}$ eigenvalue of $\mat{H}_2$.
Requires periodic boundary conditions, an FFT installation (\verb|FFTW| single and double precision on the CPU and \verb|cuFFT|/\verb|rocFFT| on GPU) and enforces that the average value of $\tildearr{\rho}$ is $0$. 

\subsection{Hypre}
The Hypre solver uses the \href{https://computing.llnl.gov/projects/hypre-scalable-linear-solvers-multigrid-methods}{Hypre library} through the \lstinline{amrex::HypreSolver}
interface to solve
\begin{equation}
    \tildearr{\rho}=-\tildemat{D}\mat{H}_2\mat{G}\arr{\phi}
\end{equation}
for $\phi$.

For more information on hypre parameters, see the \lstinline{AMReX}~\href{https://amrex-codes.github.io/amrex/docs_html/LinearSolvers.html#external-solvers}{Hypre interface documentation}.
Requires A 2D/3D CPU compilation and with periodic boundary conditions enforces that the average value of $\tildearr{\rho}$ is $0$. 

\section{Quasineutral Solver}
The quasineutral solver uses Hypre to solve the quasineutral linear system.
Requires A 2D/3D CPU compilation.
@todo: Add details here.

\end{document}